\documentclass[11pt]{article}
\usepackage[T1]{fontenc}
\usepackage{lmodern,amsmath,amssymb,booktabs,array,microtype}
\usepackage[margin=28mm]{geometry}
\usepackage[hidelinks]{hyperref}
\usepackage{listings}
\lstset{basicstyle=\small\ttfamily,breaklines=true,columns=fullflexible,
  keepspaces=true,frame=none,aboveskip=8pt,belowskip=8pt}
\setlength{\parindent}{1.2em}
\setlength{\parskip}{3pt}
\title{Formal Verification of a Structural Lower Bound\protect\\
  for Binary $3\times3$ Matrix Multiplication}
\author{Qiushi Engine Matmul Research}
\date{Companion to the research report of 9 September 2026}
\begin{document}
\maketitle
\begin{abstract}
This companion describes the Lean formalization of the structural proof that
every exact bilinear algorithm for $3\times3$ matrix multiplication over
$\mathbb F_2$ requires at least 21 products. The formal proof connects the
ordinary matrix-multiplication specification to finite quotient bounds,
occupation geometry, a forced rank profile, and an algebraic contradiction
obtained from equality in a split-flattening rank bound. Its finite premises
are proved within Lean. The
project also provides reusable results on quotient semantics, rank-additive
matrix saturation, tensor symmetry and finite subspace geometry. Sources,
certificate inputs, fixed dependencies and verification commands accompany
the mathematical argument.
\end{abstract}

\section{The Mathematical Statement}
Let $T_{333}$ be the tensor of ordinary $3\times3$ matrix multiplication over
$\mathbb F_2$. An exact bilinear algorithm computes products of a linear form
in the first input and a linear form in the second, then linearly combines
these products to recover the output. The formalization proves
\begin{equation}
  21\le R_{\mathbb F_2}(T_{333})\le23.
\end{equation}
The lower bound rules out all such algorithms using at most 20 products.
The upper bound is supplied by an
explicit 23-term decomposition, whose 729 tensor coordinates are checked.

The ordinary multiplication statement is exposed by
\texttt{bilinear\_mul\_requires\_21}. Here \texttt{Mat3} denotes
$3\times3$ matrices over \texttt{ZMod 2}:
\begin{lstlisting}
theorem bilinear_mul_requires_21 (r : Nat)
    (A B C : Fin r -> Mat3)
    (h : forall X Y : Mat3,
      X * Y = bilinearAlgorithm A B C X Y) : 21 <= r
\end{lstlisting}
This display uses ASCII equivalents for Lean notation. It quantifies over
every input matrix. Proved equivalences connect the tensor identity to the
bilinear algorithm statement, and the finite quotient bounds are established
within the proof.

\section{From Finite Bounds to Algebraic Rigidity}
Suppose a 20-term decomposition existed. Quotient-rank bounds constrain the
number of first factors in each matrix subspace. The strengthened
two-plane bounds, together with the finite geometry of rank-one differences,
force the first-factor matrix-rank profile
\begin{equation}
  (n_1,n_2,n_3)=(16,1,3).
\end{equation}
The resulting rank sum is $16+2+9=27$, exactly the rank of the full split
flattening. If $P=\sum_t M_t$ is that flattening, equality in rank
subadditivity gives
\begin{equation}
  M_tP^{-1}M_s=\delta_{ts}M_t.
\end{equation}
For the matrix-multiplication tensor, the product formula turns these
identities into constraints coupling all three tensor factors. They allow
at most one invertible first factor. The forced profile requires three,
which gives the contradiction.

Lean follows this structural argument while reorganizing intermediate
definitions and lemmas for reuse. The finite prerequisites are connected
to the same matrix space and quotient definitions as the final theorem.

\section{Finite Certificates and Coverage}
The finite proof uses exact integer branch certificates. At each leaf,
nonnegative combinations of occupation inequalities yield a contradiction;
proved integer splits cover the remaining possibilities. Transports,
quotient coordinates and the resulting rank lower bounds are proved in Lean.
The generator translates each certificate into proof terms checked by Lean.

\begin{center}
\begin{tabular}{@{}p{0.32\linewidth}p{0.60\linewidth}@{}}
\toprule
Mathematical component & Formalized result \\
\midrule
Main theorem & Lower bound 21 and ordinary bilinear semantics on all inputs \\
Finite prerequisites & Eight strengthened quotient bounds and all 496 frozen representative bounds \\
Finite geometry & Actual subspace coverage, two-plane orbits, exact labels and row bindings \\
Structural algebra & General saturation, quotient/restriction equivalence, tensor symmetry and cross-factor identities \\
Further consequences & Rank-23 witness, E11 three-term lifting implication and appropriately qualified profile restrictions \\
Encoding semantics & Integer multiplicity copies and disjoint sequential-counter CNF equivalence \\
\bottomrule
\end{tabular}
\end{center}

The precise declarations and hypotheses appear in \texttt{COVERAGE.md}.
In particular, the E11 lifting theorem is an implication from a hypothetical
core witness. Statements
about decompositions with zero summands retain the required alternatives.
The historical Python, SAT and DRAT artifacts remain separately reproducible;
the Lean proof establishes its finite premises with integer certificates
rather than formalizing the implementation or bytes of those programs.

\section{Kernel Verification}
The project pins Lean 4.33.1 and a specific Mathlib revision. All 13,438
registered modules have successful source-matched compilation outputs.
An additional six-group verification checks 89 terminal declarations and
their complete transitive dependencies using the unchanged official
\texttt{Lean.Environment.replay}, starting from an empty kernel environment.
The groups cover the main theorem and global bounds, calibration,
classification, full occupation systems, structural results and Boolean
encodings. All six checks passed.

Exact terminal types and their transitive axiom dependencies were checked
as well. The only axioms are the standard
\texttt{propext}, \texttt{Classical.choice} and \texttt{Quot.sound}.
\path{verification-results.json} records the checked inputs and outcomes,
and \path{verification.json} specifies the terminal declarations.

\section{Reproduction and Reuse}
The project runs without Qiushi Engine or language-model access. From
\texttt{formalization/}, with Elan and Python 3.11 or later installed:
\begin{lstlisting}
lake exe cache get
python3 tools/build.py . --all --jobs 4 --fresh-lean
\end{lstlisting}
Use the build directory printed by the completed command to run
\begin{lstlisting}
python3 tools/verify_all.py \
  --environment "$BUILD_DIR/environment.json" \
  --jobs 2 --timeout 43200
\end{lstlisting}
The large finite checks may take hours. The documented timeout is a ceiling,
not a guaranteed duration. Compiler and replay logs are kept locally and are
not prerequisites for another reader to rebuild the proofs.

The general algebraic interfaces can be reused separately from the full
finite computation. Quotient semantics, rank-additive saturation and tensor
symmetry provide components for extending the argument to other settings;
the certificate interfaces connect finite computations to their mathematical
consequences. The accompanying research
trajectory explains how Qiushi Engine moved from algorithm search, through
failed support-only approaches, to the equality-case argument. The written
proof explains why the result holds; the trajectory explains how the
understanding developed; Lean checks the formal chain connecting its claims.

\end{document}
